\section{Выводы}

Корпус выбран: кино и сериалы, два источника --- русская Википедия и film.ru. Они
дополняют друг друга: энциклопедические статьи и авторские рецензии. Документов хватает
с запасом, готовые поисковые системы для сравнения есть.

Главное, что я вынес из работы: \enquote{размер документа} --- не одна величина, а три.
У статьи \enquote{Во все тяжкие} это 1152~КБ сырого HTML, 161 тысяча символов текста и
18 тысяч слов, причём отношение между ними непостоянно --- от 0,4\,\% у карточки фильма
на film.ru до 23,4\,\% у интервью там же.

Ещё два наблюдения. Статьи-заглушки Википедии и карточки фильмов film.ru формально
относятся к нужной тематике, но 50--130 слов текста делают их бесполезными и занижают
все средние --- отсеивать их надо сразу. И второй источник нельзя выбирать по пяти
скачанным вручную страницам: Кино-Театр.Ру на них выглядел подходящим, а при обкачке
закрылся от робота капчей.

\subsection*{Недостатки работы}

Выборка маленькая --- 13 страниц, из них 6 по второму источнику, и отбирались они не
случайно, а специально разнотипными. Средние по такой выборке неустойчивы, тем более
при большом разбросе размеров. Метод выделения текста грубый: теряется карточка фильма,
зато в текст попадают списки источников. Объём film.ru оценён по числу страниц в
списках, а не по реально скачанным документам.

Всё это стоит пересчитать по полному корпусу, когда он будет собран, и перейти от
средних к медианам.

\pagebreak
